%-------------------------
% Julian Canales — Resume (one page)
% pdflatex JulianCanales_Resume.tex
%-------------------------
\documentclass[letterpaper,10.5pt]{article}

\usepackage[utf8]{inputenc}
\usepackage[T1]{fontenc}
\usepackage{lmodern}
\usepackage[margin=0.5in, top=0.35in, bottom=0.35in]{geometry}
\usepackage{hyperref}
\usepackage{xcolor}

\definecolor{linkblue}{HTML}{2563EB}
\hypersetup{colorlinks=true,urlcolor=linkblue,pdfauthor={Julian Canales}}

% Section heading
\renewcommand{\section}[1]{%
  \vspace{0.5em}%
  {\normalsize\bfseries\scshape #1}%
  \vspace{-0.5em}%
  \par\noindent\rule{\textwidth}{0.4pt}%
  \vspace{0.15em}%
}

% Tight itemize
\let\olditemize\itemize
\let\oldenditemize\enditemize
\renewenvironment{itemize}{%
  \olditemize
  \setlength{\itemsep}{1pt}%
  \setlength{\parskip}{0pt}%
  \setlength{\parsep}{0pt}%
  \setlength{\topsep}{1.5pt}%
  \setlength{\leftmargin}{1.2em}%
}{\oldenditemize}

\pagestyle{empty}
\setlength{\parindent}{0pt}

\begin{document}

%----------HEADER----------
\begin{center}
  {\large\bfseries Julian Canales} \\[2pt]
  {\small 832-931-1132 \enspace$\vert$\enspace
  \href{mailto:juliancanales@utexas.edu}{juliancanales@utexas.edu} \enspace$\vert$\enspace
  San Francisco, CA \enspace$\vert$\enspace
  \href{https://github.com/juliancanaless}{GitHub} \enspace$\vert$\enspace
  \href{https://www.linkedin.com/in/juliancanales05/}{LinkedIn}}
\end{center}

%----------EDUCATION----------
\section{Education}

{\small
\textbf{The University of Texas at Austin} --- M.S.\ Artificial Intelligence (online) \hfill Aug 2025 -- May 2027 \\
Coursework: Machine Learning, Online Learning \& Optimization, Deep Learning

\vspace{3pt}
\textbf{The University of Texas at Austin} --- B.S.\ Computer Science \& Mathematics \hfill May 2025 \enspace$\vert$\enspace GPA: 3.65 \\
\textit{Awards:} QuestBridge Scholar, HSF Scholar, University Honors \\
\textit{Activities:} Directed Reading Program, Texas Blockchain, CS Internal Transfer Program \\
\textit{Coursework:} Principles of ML, Distributed Computing, Concurrency (Honors), Geometric Foundations of Data Science, Quantum Info Sci, Real Analysis I--II, Stochastic Processes I (Grad), Predictive Analytics, Game Theory
}

%----------EXPERIENCE----------
\section{Experience}

{\small
\textbf{Avathon} \enspace$\vert$\enspace Data Scientist \hfill Oct 2025 -- Present
\begin{itemize}
  \item Built AI-oriented microservices and agent workflows for internal infrastructure, focusing on reliability and extensibility.
  \item Built an AI agent fluent in the knowledge graph's query language (IQL) with tools for schema discovery, graph traversal, and embeddings-backed search; required deep research into the CKG codebase to surface undocumented capabilities and validate embedding pipelines.
  \item Worked full-stack on a transport management system, including resource-selection config and operational logic.
  \item Pitched and helped adopt DBOS for workflow durability and orchestration across services.
  \item Integrated Langfuse for AI observability; continuing work on AI infrastructure, frontend architecture, and DevOps.
\end{itemize}

\textbf{IBM} \enspace$\vert$\enspace Data Science Intern \hfill May -- Aug 2023
\begin{itemize}
  \item Built and deployed a Flask POC integrating Watson Assistant with Db2; helped close a client deal.
  \item Created a 1.5k-LOC Dash app backed by watsonx, Discovery, and a custom eval API for LLM comparisons.
  \item Owned backend API integration for a React + Flask onboarding MVP on IBM Cloud and OpenShift.
\end{itemize}

\textbf{UT Austin} \enspace$\vert$\enspace Undergraduate Researcher, Earthquake Modeling \hfill Jan -- May 2024
\begin{itemize}
  \item Re-engineered a legacy Mathematica SDOF solver into a multithreaded Python pipeline.
  \item Processed 29k+ NGA-West2 records; ran 64-vCPU GCP sweeps, cutting runtime from $\approx$3\,yrs to $\approx$2\,wks.
  \item Produced interactive Plotly dashboards and a seminar deck for faculty presentation.
\end{itemize}
}

%----------PROJECTS----------
\section{Selected Projects}

{\small
\textbf{Crabs AI Assistant} \enspace$\vert$\enspace \textit{OpenClaw, LightRAG, CouchDB} \hfill 2026
\begin{itemize}
  \item Self-hosted AI assistant running on a Hetzner VPS with CouchDB syncing an Obsidian vault across devices.
  \item Two-tiered memory system: short-term conversational recall via OpenClaw and long-term knowledge graph traversal via LightRAG; zero-guessing policy forces the agent to always search actual notes instead of relying on pre-trained weights.
\end{itemize}

\textbf{Adaptive DQN Planner} \enspace$\vert$\enspace \textit{RL, DQN, CEM, Python} \hfill 2025
\begin{itemize}
  \item Built an RL agent for autonomous driving in a roundabout simulator, pairing DQN with a CEM optimizer and k-NN error tracker to adapt uncertainty estimates during planning.
  \item Achieved 0\% collisions and +17\% safety margin over deterministic baselines across 20 evaluation runs.
\end{itemize}

\textbf{CUDA K-Means / K-Means++} \enspace$\vert$\enspace \textit{CUDA, C++} \hfill 2025
\begin{itemize}
  \item Implemented 4 GPU-accelerated K-Means variants including shared memory and K-Means++ initialization; benchmarked on dual Quadro RTX 6000s, achieving 14x speedup on 1M$\times$32D datasets.
\end{itemize}

\textbf{GPU Fine-Grained Sync for EM} \enspace$\vert$\enspace \textit{CUDA, Distributed Systems} \hfill 2025
\begin{itemize}
  \item Recreated and extended Wang et al.\ (ASPLOS '19) fine-grained GPU synchronization using a client-server design with shared-memory buffers; benchmarked on counters, hash tables, and EM for GMMs.
\end{itemize}

\textbf{Fantasy Draft RL Agent} \enspace$\vert$\enspace \textit{Python, PyTorch, Gymnasium} \hfill 2024
\begin{itemize}
  \item Built a custom Gymnasium environment modeling 12-team snake drafts with 300+ players, action masking, and roster-aware observations; trained Maskable-PPO across multiple NFL seasons, outperforming heuristic baselines by 26--42\%.
\end{itemize}

\textbf{Black-Scholes Derivation Paper} \enspace$\vert$\enspace \textit{Probability, Measure Theory, Stochastic Calculus} \hfill 2024
\begin{itemize}
  \item Wrote a self-contained derivation of Black-Scholes for European call options from first principles, covering measure theory, martingales, Brownian motion, and connecting the PDE framework to practical asset pricing.
\end{itemize}
}

%----------SKILLS----------
\section{Skills}

{\small
\textbf{Languages:} Python, SQL, C/C++, CUDA C, TypeScript, Bash \enspace$\vert$\enspace \textit{Familiar:} Java, Go, Rust, R \\
\textbf{AI \& Agents:} LLM APIs (OpenAI, Anthropic, Google), agent orchestration \& tooling, prompt engineering, RAG pipelines, LiteLLM, ChromaDB, LightRAG, Langfuse, model routing \& compatibility, streaming/thinking pipelines \\
\textbf{ML \& Data:} PyTorch, scikit-learn, pandas, NumPy, SciPy, Gymnasium, Stable-Baselines3, Matplotlib/Plotly \\
\textbf{Web \& APIs:} Flask, FastAPI, Django, React, Next.js, REST APIs, Dash, Tailwind CSS, SSE/WebSockets \\
\textbf{Infra \& Tools:} AWS (EC2/S3), GCP, Docker, Git, Kubernetes, OpenMP/MPI, GPU profiling (nvprof/nsight) \\
\textbf{Spoken:} English (native), Spanish (native)
}

\end{document}
